\chapter{Introduction}\label{ch:introduction}

\ps{}~\cite{pystencils} describes itself as \emph{``a symbolic domain-specific language
and metaprogramming toolkit for writing high-performance numerical stencil kernels for a
variety of hardware targets''}. This manual explains what that means in practice. It
covers the problem \ps{} solves, how you use it as a programmer, what goes in and what
comes out, and how the code generator works inside, including its backends.

\section{What is a stencil code?}

Many numerical methods update every cell of a regular grid using only the values of a
few neighbouring cells. The fixed pattern of neighbours is called a \term{stencil}. A
typical example is one step of a Jacobi solver for the Laplace equation on a 2D grid:
\begin{equation}
  u^{\text{new}}_{i,j} \;=\; \tfrac14\,\bigl(u_{i+1,j} + u_{i-1,j} + u_{i,j+1} + u_{i,j-1}\bigr).
  \label{eq:jacobi}
\end{equation}
The same pattern appears in finite-difference and finite-volume schemes, in image
filters, in cellular automata and in lattice Boltzmann methods~\cite{lbmpy}. A
\term{stencil kernel} is the function that applies such an update to every cell of an
array. Cells at the border do not have all their neighbours, so the arrays usually carry
extra \term{ghost layers} around the domain that the kernel reads but does not update.
\Cref{fig:stencil} shows one application of \cref{eq:jacobi}: four values of the
source array are read and one value of the destination array is written. A kernel
repeats this for every interior cell.

\begin{figure}[htbp]
  \centering
  % One application of the 2D Jacobi stencil: read 4 neighbours in src, write 1 cell in dst.
\begin{tikzpicture}[x=0.62cm,y=0.62cm,font=\footnotesize]
  % ---- src -------------------------------------------------------------------
  \fill[ghostcell] (0,0) rectangle (7,7);
  \fill[white] (1,1) rectangle (6,6);
  \cellfill[nbcell]{4}{3} \cellfill[nbcell]{2}{3}
  \cellfill[nbcell]{3}{4} \cellfill[nbcell]{3}{2}
  \fill[pattern=north east lines,pattern color=psblue!60] (3,3) rectangle ++(1,1);
  \draw[gridline] (0,0) grid (7,7);
  \draw[domainline] (1,1) rectangle (6,6);
  \node[celltext] at (4.5,3.5) {E};  \node[celltext] at (2.5,3.5) {W};
  \node[celltext] at (3.5,4.5) {N};  \node[celltext] at (3.5,2.5) {S};
  \node[above,font=\small\bfseries] at (3.5,7) {\texttt{src}};
  % callouts
  \draw[psgray] (4,4.5) -- (7.4,4.5) node[right,celltext] {\texttt{src[0,1]}\,(N)};
  \draw[psgray] (5,3.5) -- (7.4,3.5) node[right,celltext] {\texttt{src[1,0]}\,(E)};
  \draw[psgray] (4,2.5) -- (7.4,2.5) node[right,celltext] {\texttt{src[0,-1]}\,(S)};
  \draw[psgray] (2,3.5) -- (-0.4,3.5) node[left,celltext] {(W)\,\texttt{src[-1,0]}};
  % axes
  \draw[axisarrow] (0,-0.6) -- (2,-0.6) node[right,celltext] {$x$ (\texttt{ctr\_0})};
  \draw[axisarrow] (-0.6,0) -- (-0.6,2) node[above,celltext] {$y$ (\texttt{ctr\_1})};

  % ---- dst -------------------------------------------------------------------
  \begin{scope}[shift={(12.5,0)}]
    \fill[ghostcell] (0,0) rectangle (7,7);
    \fill[white] (1,1) rectangle (6,6);
    \cellfill[centercell]{3}{3}
    \draw[gridline] (0,0) grid (7,7);
    \draw[domainline] (1,1) rectangle (6,6);
    \node[above,font=\small\bfseries] at (3.5,7) {\texttt{dst}};
    \draw[psgray] (3.5,3) -- (3.5,-0.5) node[below,celltext] {\texttt{dst[0,0]} = \texttt{dst.center()}};
  \end{scope}

  % ---- the update ------------------------------------------------------------
  \draw[flowarrow,psblue] (5.6,6.3) to[bend left=22]
    node[above,align=center,font=\footnotesize,text=black]
      {$\tfrac14\,(\text{E}+\text{W}+\text{N}+\text{S})$}
    (15.6,3.9);

  % ---- legend ----------------------------------------------------------------
  \begin{scope}[shift={(0,-2.4)}]
    \fill[ghostcell,draw=psgray!55] (0,0) rectangle ++(0.7,0.7);
    \node[right,celltext] at (0.8,0.35) {ghost layer (only read)};
    \fill[nbcell,draw=psgray!55] (7.6,0) rectangle ++(0.7,0.7);
    \node[right,celltext] at (8.4,0.35) {cells read};
    \fill[centercell,draw=psgray!55] (12.4,0) rectangle ++(0.7,0.7);
    \node[right,celltext] at (13.2,0.35) {cell written};
    \draw[domainline] (17.4,0) rectangle ++(0.7,0.7);
    \node[right,celltext] at (18.2,0.35) {domain};
  \end{scope}
\end{tikzpicture}

  \caption[One application of the Jacobi stencil]{One application of the Jacobi stencil of \cref{eq:jacobi}: the four neighbours (orange) of a cell are read from \texttt{src} and their mean is written to the same cell (blue) of \texttt{dst}. The hatched cell is where the result belongs; its own value is not used. The grey ghost layer is read but never written.}
  \label{fig:stencil}
\end{figure}

Stencil kernels are simple to state, but fast implementations are not:
\begin{itemize}
  \item The loop nest must follow the array's \emph{memory layout} (which index is
        contiguous?), or performance can drop by an order of magnitude.
  \item Index arithmetic with \emph{strides} and ghost layers is easy to get wrong,
        especially for multi-component fields (e.g.\ 19 or 27 values per cell).
  \item Good CPU performance needs \emph{SIMD vectorization} (SSE/AVX/AVX-512, NEON,
        SVE, RISC-V V) and \emph{multithreading} (OpenMP). Each SIMD instruction set
        has its own intrinsics.
  \item GPUs need a mapping from cells to threads and blocks, a launch configuration,
        and different math functions.
  \item Real models contain hundreds of floating-point operations per cell. Shared
        subexpressions must be computed only once, in the right precision.
\end{itemize}
Writing and maintaining one hand-tuned variant per target, precision and layout does
not scale. Each variant has to be kept consistent with the mathematical model.

\section{The pystencils idea: describe, then generate}

\ps{} separates \emph{what} a kernel computes from \emph{how} it is executed:
\begin{enumerate}
  \item You describe the update rule \textbf{symbolically} in Python with
        SymPy~\cite{sympy}: \term{fields} (abstract $n$-dimensional arrays) and
        \term{assignments} such as \code{dst[0,0] = (src[1,0] + src[-1,0] + ...) / 4}.
        At this level the description contains no loops, no strides, no types and no
        intrinsics.
  \item You choose a \textbf{configuration}: the hardware target, the data types, the
        iteration space and optimizations such as OpenMP or vectorization.
  \item The \textbf{code generator} \code{ps.create_kernel} turns description and
        configuration into one optimized C/C++ function. It handles index arithmetic,
        the loop order, peeling and remainder loops, intrinsics, the GPU thread mapping
        and type checking.
  \item The resulting \textbf{kernel object} can be \emph{just-in-time compiled} and
        called from Python with NumPy (or CuPy) arrays. Its code and metadata can also be
        \emph{exported} into a C++ application such as waLBerla~\cite{walberla}.
\end{enumerate}
Because the model is a symbolic expression, it can be transformed \emph{before} code
generation. You can discretize partial differential equations automatically
(\code{ps.fd}), eliminate common subexpressions (\code{ps.simp}), or derive whole
families of models, as lbmpy~\cite{lbmpy} does for lattice Boltzmann methods.

\begin{keyidea}[Key idea]
  A stencil is \emph{data}, a SymPy expression, rather than code. \ps{} is a
  \emph{compiler} from that data to target-specific C/C++. One symbolic description
  yields scalar, SIMD, OpenMP, CUDA, HIP and SYCL kernels in single or double precision.
\end{keyidea}

\section{Inputs and outputs at a glance}

\Cref{fig:big-picture} summarizes the data flow covered in this manual.
\Cref{tab:io} lists the most important objects on each side.

\begin{figure}[htbp]
  \centering
  % Big picture: inputs -> pystencils (create_kernel) -> Kernel object -> its uses
\begin{tikzpicture}[
  font=\small,
  box/.style={draw=#1!70!black,fill=#1!8,rounded corners=2pt,align=left,
    inner sep=5pt,text width=3.75cm},
  core/.style={draw=psblue!80!black,fill=psblue!7,rounded corners=2pt,align=center,
    inner sep=4pt,text width=3.3cm,minimum height=0.75cm},
  arr/.style={-{Stealth[length=2.2mm]},thick,draw=psgray},
  lbl/.style={font=\footnotesize\bfseries,text=psgray},
]
  % inputs
  \node[box=psteal] (in1) {\textbf{Fields}\\[1pt]\footnotesize\texttt{ps.fields("f, g: double[2D]")}};
  \node[box=psteal,below=0.25cm of in1] (in2) {\textbf{Assignments} (SymPy)\\[1pt]\footnotesize\texttt{g[0,0] = (f[1,0]+f[-1,0])/2}};
  \node[box=psteal,below=0.25cm of in2] (in3) {\textbf{Configuration}\\[1pt]\footnotesize target, data types, iteration space, OpenMP, SIMD, GPU indexing};
  \node[lbl,above=0.15cm of in1] {INPUT (Python)};

  % core
  \node[core,right=1.2cm of in1.north east,anchor=north west] (c1) {Analysis and checks\\\footnotesize\texttt{KernelAnalysis}};
  \node[core,below=0.22cm of c1] (c2) {Translation to typed IR\\\footnotesize freeze + typify};
  \node[core,below=0.22cm of c2] (c3) {Transformations\\\footnotesize loops, SIMD, OpenMP, GPU};
  \node[core,below=0.22cm of c3] (c4) {Platform lowering\\\footnotesize intrinsics, math functions};
  \node[core,below=0.22cm of c4] (c5) {Finalize\\\footnotesize parameters, headers};
  \draw[arr] (c1) -- (c2); \draw[arr] (c2) -- (c3); \draw[arr] (c3) -- (c4); \draw[arr] (c4) -- (c5);
  \begin{scope}[on background layer]
    \node[draw=psblue!60,dashed,rounded corners=3pt,fit=(c1)(c5),inner sep=6pt,
      label={[lbl]above:\ps\ \texttt{create\_kernel}}] (core) {};
  \end{scope}

  % outputs: the Kernel object and what is done with it
  \node[box=psorange,anchor=north west] (o1) at ($(core.east |- c1.north)+(1.2,0)$)
    {\textbf{C/C++ source code}\\[1pt]\footnotesize \texttt{kernel.get\_c\_code()}: plain C,
     intrinsics, CUDA/HIP, SYCL};
  \node[box=psorange,below=0.25cm of o1] (o2)
    {\textbf{Python callable} (JIT)\\[1pt]\footnotesize \texttt{kernel.compile()}, call
     with NumPy / CuPy arrays};
  \node[box=psorange,below=0.25cm of o2] (o3)
    {\textbf{Embedded in C++}\\[1pt]\footnotesize \texttt{no\_jit} + pystencils-sfg,
     waLBerla};
  \node[lbl,above=0.15cm of o1] {OUTPUT};
  \node[box=psorange,below=0.4cm of o3,line width=1pt] (k)
    {\textbf{\texttt{Kernel} object}\\[1pt]\footnotesize body (IR), parameters with
     properties, headers, metadata};

  \draw[arr] (in1.east) -- (c1.west);
  \draw[arr] (in2.east) -- ++(0.6,0) |- (c1.west);
  \draw[arr] (in3.east) -- ++(0.6,0) |- (c1.west);
  \draw[arr] (c5.east) -- ++(0.55,0) |- (k.west);
  \foreach \o in {o1,o2,o3} \draw[arr] (k.east) -- ++(0.3,0) |- (\o.east);
\end{tikzpicture}

  \caption[From a symbolic description to executable kernels]{From a symbolic
    description to executable kernels. The code generator (middle) produces a
    \texttt{Kernel} object; the C/C++ code, the compiled Python function and the
    integration into C++ applications are all derived from it. \Cref{ch:architecture}
    looks at the code generator in detail.}
  \label{fig:big-picture}
\end{figure}

\begin{table}[htbp]
  \centering\small
  \caption{The main user-facing objects.}\label{tab:io}
  \begin{tabularx}{\linewidth}{@{}l>{\raggedright\arraybackslash}X l@{}}
    \toprule
    \textbf{Object} & \textbf{Role} & \textbf{Chapter}\\
    \midrule
    \code{ps.Field}, \code{ps.fields(...)} & abstract arrays and neighbour accesses \code{f[1,0]} & \ref{ch:symbolic-frontend}\\
    \code{ps.Assignment}, \code{ps.AssignmentCollection} & the update rule, optionally with subexpressions & \ref{ch:symbolic-frontend}\\
    \code{ps.TypedSymbol}, \code{ps.tcast} & scalars with explicit types; casts & \ref{ch:symbolic-frontend}\\
    \code{ps.CreateKernelConfig}, \code{ps.Target} & target, types, iteration space, optimizations & \ref{ch:kernel-creation}\\
    \code{ps.create_kernel(...)} & the code generator entry point & \ref{ch:kernel-creation}\\
    \midrule
    \code{Kernel} / \code{GpuKernel} & generated function: IR, parameters, headers, metadata & \ref{ch:outputs-execution}\\
    \code{kernel.get_c_code()} & the C/C++ source text & \ref{ch:outputs-execution}\\
    \code{kernel.compile()} & a Python callable (\code{KernelWrapper}) & \ref{ch:outputs-execution}\\
    \bottomrule
  \end{tabularx}
\end{table}

\section{The pystencils ecosystem}

\ps{} is developed at Friedrich-Alexander-Universität Erlangen-Nürnberg as part of the
\emph{pycodegen} family of tools:
\begin{description}
  \item[lbmpy~\cite{lbmpy}] builds lattice Boltzmann methods (collision operators,
        equilibria, boundary conditions) symbolically. It hands the resulting assignment
        collections to \ps{} for code generation.
  \item[waLBerla~\cite{walberla}] is a massively parallel C++ multiphysics framework.
        Its code generation layer (the newer \emph{sweepgen}, built on
        pystencils-sfg) embeds \ps{} kernels into C++ \emph{sweeps}. These run on
        block-structured, distributed grids.
  \item[pystencils-sfg~\cite{pystencils-sfg}] (\emph{source file generator}) writes
        complete C++ source and header files around \ps{} kernels. It maps C++ data
        structures to the kernel parameters, as described in
        \cref{sec:external-use}.
\end{description}

\section{Scope of this manual}

The manual is based on the \ps{} sources in the \file{submodules/pystencils} submodule
of this workspace. That is version \texttt{2.0b1+4.g9a989e16} of the \emph{pystencils
2.0} line, a near-complete rewrite of the 1.x series. The legacy 1.4 releases live on the
\texttt{v1.x} branch upstream. The submodule is checked out on the branch
\texttt{rvv-support}, which adds a RISC-V vector backend on top of upstream \ps. That
backend is covered in \cref{ch:rvv-backend} and marked with boxes like this one:

\begin{forkonly}
  Material in an orange box describes code that exists only on the \texttt{rvv-support}
  branch of this workspace, not in upstream \ps.
\end{forkonly}

\paragraph{How to read this manual.}
\Cref{ch:getting-started,ch:symbolic-frontend,ch:kernel-creation,ch:outputs-execution}
(Part~I) are written for \emph{users}. They teach the workflow from fields to running
kernels. \Cref{ch:architecture,ch:ir,ch:translation,ch:transformations,ch:platforms,ch:rvv-backend}
(Part~II) are for \emph{developers}. They explain the internal representation, the
compiler passes and the hardware platforms. \Cref{ch:worked-example} closes with a
complete lattice Boltzmann example. The appendices contain a quick reference
(\cref{app:quick-reference}), a reference of all configuration options, a map of the
source tree, a list of pitfalls, and a glossary.

\paragraph{The listings are real.}
The example code in this manual comes from the scripts in
\file{docs/manual_pystencil/examples/}. All generated code, IR dumps and printed values
were captured by running those scripts against the submodule with
\file{./build.sh listings}. The scripts check their own results, mostly against NumPy.
Generated C code was re-indented with \texttt{clang-format} to fit the page. Only
whitespace and line breaks differ from what \ps{} prints. The only exceptions are a few
short fragments that show a signature, a configuration pattern or an excerpt of the
\ps{} sources. They are typed directly into the text and not captured from a run.
